\chapter{實驗成果分析與驗證}
\label{chapter:experiments}

第三章建立了 PI-ST-GNN 之完整方法論架構。本章之目的並非展示運算速度的極致最佳化，而是回答一個對建築資訊學實務更關鍵的問題：這套物理資訊圖神經網路，是否真正學到了都市微氣候的空間因果結構，其預測結果是否足以支撐建築師的設計決策。本章依序呈現訓練配置、測試集量化驗證、消融實驗之設計選擇分析，以及運算效能對互動式設計工作流的實質意義。


\section{訓練流程與超參數配置}
\label{sec:training}

\subsection{訓練環境與硬體配置}
\label{subsec:hardware}

本研究所有訓練與推理實驗均於搭載下列硬體之工作站執行：

\begin{itemize}
    \item \textbf{GPU：}NVIDIA GeForce RTX 5070 Ti Laptop GPU（視訊記憶體 12.8 GB）
    \item \textbf{CPU：}AMD Ryzen AI 9 HX（架構：Zen 5）
    \item \textbf{框架：}PyTorch 2.x + 自製輕量圖引擎（無 PyG 依賴）
\end{itemize}


\subsection{訓練超參數與學習率排程}
\label{subsec:training_params}

PI-ST-GNN 之訓練遵循下列配置，詳見表~\ref{tab:training_config}：

\begin{table}[htbp]
    \centering
    \caption{模型訓練超參數配置}
    \label{tab:training_config}
    \small
    \begin{tabular}{llc}
        \toprule
        \textbf{超參數} & \textbf{說明} & \textbf{設定值} \\
        \midrule
        批次大小（Batch Size）         & 每批場景數（逐場景）            & 1 \\
        最大訓練週期（Max Epochs）     & 允許最長訓練回合數              & 250 \\
        早停耐心（Early Stopping）     & 驗證損失無改善之容忍週期數      & 20 \\
        初始學習率（LR）               & Adam 最佳化器初始步長             & $1 \times 10^{-3}$ \\
        學習率排程（LR Scheduler）     & ReduceLROnPlateau，觸發後對半   & factor = 0.5 \\
        學習率下限（Min LR）           & 排程器觸發之最小學習率下限      & $\approx 7.8 \times 10^{-6}$ \\
        資料損失函數（$\mathcal{L}_{\text{data}}$）& 預測值與標籤之 MSE  & --- \\
        物理損失權重（$\alpha_1, \alpha_2, \alpha_3$）& 輻射、時序、風場    & 各自調校 \\
        Dropout                        & RGCN 訓練期間 Dropout 率        & 0.1 \\
        \bottomrule
    \end{tabular}
\end{table}

學習率排程採用 ReduceLROnPlateau 策略：當驗證損失在連續數個 Epoch 未見改善時，學習率縮減為原來的 0.5 倍。本研究最終採用之 V5-300 模型（真實場景資料集，Radiance／EnergyPlus 模擬引擎，見第~\ref{sec:v4_validation}節）訓練中，學習率於第 23、111、130 Epoch 觸發三次衰減（自 $10^{-3}$ 逐步降至 $1.25 \times 10^{-4}$），模型於第 139 Epoch 因連續 20 個 Epoch 無改善觸發早停（Early Stopping）而結束訓練，最佳權重取自第 \textbf{119} Epoch，驗證 $R^2$ 達到 \textbf{0.9837}。

\subsection{訓練曲線與收斂特性}
\label{subsec:convergence}

圖~\ref{fig:training_curves} 展示訓練損失、驗證損失與驗證 $R^2$ 之逐 Epoch 演化曲線，其分段特性揭示模型學習都市微氣候映射關係的三個階段：初期（前約 10 個 Epoch）驗證 $R^2$ 自暖身階段急速上升，顯示模型迅速掌握異質圖節點特徵與熱壓力間的基本映射；中期（Epoch 10--100）$R^2$ 持續波動精調，對應真實建物遮蔽、植栽蒸散、對流與真實地表熱傳導等物理管道之關係權重逐步分化學習；後期（Epoch 100 後）訓練與驗證損失同步趨於平穩收斂，驗證 $R^2$ 穩定於 0.98 附近，未見發散，確認物理資訊損失函數對解空間提供了有效的正則化，模型未發生過擬合。

\begin{figure}[htbp]
    \centering
    \includegraphics[width=\textwidth]{../urban-thermal-gnn/04_training/viz_output/training_v5_300/figA_training_curves.png}
    \caption[模型訓練收斂曲線]{模型訓練收斂曲線（V5-300，共 139 Epoch，早停於最佳 Epoch 119）}
    \label{fig:training_curves}
\end{figure}

上圖上半部為訓練損失（藍）與驗證損失（橘）之逐 Epoch MSE 演化，灰色垂直線標示學習率衰減時間點；下半部為驗證 $R^2$ 逐 Epoch 提升曲線，最終收斂至最佳 $R^2 = 0.9837$（Epoch 119，紅點標示）。


\section{測試集量化評估結果}
\label{sec:evaluation}

\subsection{整體預測精確度指標}
\label{subsec:overall_metrics}

在 45 個真實場景測試集、共計 166,188 個預測點次（各場景空氣節點數依真實建物量體與場地邊界而異，介於 136--400 節點間，逐場景 $\times$ 11 時步加總）上，PI-ST-GNN（V5-300：真實 OSM／ETH／IoT 地理資料 $+$ Radiance／EnergyPlus 模擬引擎，見第~\ref{sec:v4_validation}節）之最終測試集評估結果如表~\ref{tab:main_results} 所示。

\begin{table}[htbp]
    \centering
    \caption[測試集評估結果]{PI-ST-GNN（V5-300）測試集評估結果（共 166,188 預測點次，45 個真實測試場景）}
    \label{tab:main_results}
    \small
    \begin{tabular}{lcccc}
        \toprule
        \textbf{指標} & \textbf{$R^2$} & \textbf{RMSE（°C）} & \textbf{MAE（°C）} & \textbf{等級準確率} \\
        \midrule
        PI-ST-GNN（本研究） & \textbf{0.9753} & \textbf{0.833} & \textbf{0.524} & \textbf{93.9\%} \\
        \bottomrule
    \end{tabular}
\end{table}

四項指標對應四個不同的設計輔助情境，缺一不可：決定係數 $R^2 = 0.9753$ 代表模型能解釋 UTCI 變異量的 97.53\%，超過一般代理模型部署所需之 $R^2 > 0.97$ 門檻 \citep{westermann2019surrogate}；均方根誤差 RMSE $= 0.833\,^\circ\text{C}$（去正規化後之實際溫度單位）遠小於 UTCI 熱壓力等級間距（最小間距 $6\,^\circ\text{C}$），確保等級判斷之邊界不因模型誤差而錯置 \citep{fiala2011utci}；平均絕對誤差 MAE $= 0.524\,^\circ\text{C}$ 低於 RMSE，反映多數評估點具有極高逐點精確度、少數離群值未顯著拉高整體偏差；熱壓力等級分類準確率 93.9\% 則是實際設計輔助情境下最直接相關的指標——建築師關心的通常不是溫度小數點後第二位，而是某一開放空間是否落在「強烈」或「極端」熱壓力等級。此組數字係於 Radiance／EnergyPlus 模擬之真實場景上取得，較純程序化幾何之數值指標略低，其原因與方法論意義詳見第~\ref{sec:v4_validation}節。

\begin{figure}[htbp]
    \centering
    \includegraphics[width=0.85\textwidth]{../urban-thermal-gnn/04_training/viz_output/training_v5_300/figB_scatter_test.png}
    \caption[測試集預測值散點圖]{測試集預測值與模擬標籤之散點圖（V5-300，共 166,188 個樣本）}
    \label{fig:scatter_test}
\end{figure}

上圖點雲緊密分佈於對角線 $y=x$ 附近，色彩代表密度（黃色為高密度區域），圖中所示分布形態與表~\ref{tab:main_results}之 $R^2=0.9753$ 一致，顯示模型在全 UTCI 值域內均具優秀的預測一致性。


\subsection{逐時段預測精確度分析}
\label{subsec:hourly_r2}

圖~\ref{fig:hourly_r2} 展示 11 個時步（08:00--18:00）之個別 $R^2$ 值。所有時步之 $R^2$ 均維持在良好水準，充分驗證模型對日間熱壓力動態的全面捕捉能力；高峰熱壓力時段（10:00--14:00）因太陽高度角最大、陰影梯度最為顯著，RGCN 的空間幾何建模能力得以充分發揮，$R^2$ 表現因而相對穩定。此一時段性差異本身即是模型確實學到「陰影邊界＝關鍵訊號」此一物理直覺的間接證據，而非單純記憶訓練分佈。

\begin{figure}[htbp]
    \centering
    \includegraphics[width=0.8\textwidth]{../urban-thermal-gnn/04_training/viz_output/training_v5_300/figC_hourly_r2.png}
    \caption[逐時段驗證 $R^2$ 分佈]{11 個時步（08:00--18:00）之逐時段驗證 $R^2$ 分佈（V5-300）}
    \label{fig:hourly_r2}
\end{figure}


\subsection{熱壓力等級分類混淆矩陣}
\label{subsec:confusion_matrix}

圖~\ref{fig:confusion} 以 UTCI 熱壓力六等分類體系之混淆矩陣呈現分類品質。高頻等級（中等熱壓力 26--32°C、強烈熱壓力 32--38°C）之分類準確率最高，反映模型在訓練資料密集區域的優異表現；較稀有等級（極端熱壓力 $>46°C$）之準確率相對略低，源於真實場景中該等級樣本數天然較少，而非模型結構性缺陷；絕大多數誤分案例均發生於相鄰等級之臨界點，不存在跨兩個等級以上的嚴重誤判，顯示模型誤差具連續性而非系統性偏移，此點對於建築師解讀預測結果的可信度至關重要。

\begin{figure}[htbp]
    \centering
    \includegraphics[width=0.65\textwidth]{../urban-thermal-gnn/04_training/viz_output/training_v5_300/figD_confusion_matrix.png}
    \caption[熱壓力等級分類混淆矩陣]{測試集 UTCI 熱壓力六等級分類混淆矩陣（V5-300，正規化百分比）}
    \label{fig:confusion}
\end{figure}

上圖對角線元素代表各等級之正確分類率；非對角線元素代表鄰等級間之混淆。各等級準確率分別為：中等熱壓力 88.97\%、強烈熱壓力 92.90\%、極強熱壓力 95.35\%、極端熱壓力 91.50\%。


\subsection{空間熱壓力場視覺化比較}
\label{subsec:spatial_comparison}

圖~\ref{fig:spatial_comparison} 以選定測試場景展示模型預測之空間 UTCI 熱壓力場與 Radiance／EnergyPlus 模擬標籤的視覺比較。預測場與標籤場在建築陰影輪廓、高熱區域分佈及梯度走向上均高度吻合，殘差集中於建築邊緣之陰影過渡帶——此一誤差分布特性本身具設計意義：模型在「陰影是否存在」的判斷上高度可靠，誤差主要來自陰影邊界的精確位置這類次要幾何細節，對建築師進行量體配置之初期決策影響有限。

\begin{figure}[htbp]
    \centering
    \includegraphics[width=\textwidth]{../urban-thermal-gnn/04_training/viz_output/training_v5_300/figE_spatial_comparison.png}
    \caption[空間 UTCI 熱壓力場比較]{測試場景空間 UTCI 熱壓力場比較（V5-300，正午 12:00 時步）}
    \label{fig:spatial_comparison}
\end{figure}

上圖左列為物理式模擬標籤，右列為 PI-ST-GNN 預測值，下方色條代表 UTCI 值域（°C），殘差圖（右下）顯示誤差主要集中在建築邊緣之陰影過渡帶。

圖~\ref{fig:feature_fields} 進一步展示輸入特徵場與 UTCI 輸出場之空間對應關係，並沿時間軸呈現其逐時演變：天空視角因子（SVF）對 UTCI 分佈具顯著驅動作用，高 SVF（開放地面）對應高 UTCI、低 SVF（建築陰影）對應低 UTCI，此一單調對應關係之忠實還原，正是第~\ref{sec:physics_loss}節輻射一致性約束發揮作用的直接證據；而此一對應關係隨時間仍保持穩定（各列間色彩分布之相對格局一致），顯示模型並未僅學到單一時步之表面統計巧合。

\begin{figure}[htbp]
    \centering
    \includegraphics[width=\textwidth]{../urban-thermal-gnn/04_training/figures/fig_geo6_feature_fields.png}
    \caption[輸入特徵場與輸出場空間分佈]{真實場景（V5）輸入特徵場與 UTCI 輸出場空間分佈：四個代表時步（08:00、11:00、14:00、18:00）對照}
    \label{fig:feature_fields}
\end{figure}

上圖由左至右分別為天空視角因子（SVF，靜態）、空氣溫度（Ta）、平均輻射溫度（MRT）及模型預測 UTCI 值，由上至下為四個代表時步；灰色多邊形為真實建物量體，黑點標示該時步位於陰影中之空氣節點。


\section{設計選擇之消融實驗}
\label{sec:ablation}

高整體精度本身不足以說明模型「為何」有效——若不拆解各設計選擇之獨立貢獻，便無從判斷 RGCN 的異質圖結構、抑或物理引導損失，何者才是模型可信賴性的真正來源。本節以五組消融變體（A0--A6，編號沿用各變體所加入之核心模組層次，與第~\ref{sec:v4_validation}節之資料集版本編號 V1--V5 為互不相關之獨立編號體系）系統性地移除或加入各核心模組，於真實場景資料集（V5-300 相同之訓練／驗證／測試切分）上評估其獨立貢獻。

\subsection{消融實驗設計}
\label{subsec:ablation_design}

五個消融變體之設計邏輯如下：

\begin{itemize}
    \item \textbf{A0（MLP 基準線）：}以純多層感知器（MLP）替換整個 GNN 架構，直接從節點座標與氣象統計特徵映射至 UTCI 輸出，不具空間圖建模能力，作為「無空間感知」基準。
    \item \textbf{A3（GNN 無物理損失）：}保留完整 RGCN + LSTM 架構，但訓練時移除所有物理資訊損失項（$\alpha_1 = \alpha_2 = \alpha_3 = 0$），僅以資料損失 $\mathcal{L}_{\text{data}}$ 訓練。
    \item \textbf{A4（GNN + $\mathcal{L}_{\text{rad}}$）：}在 A3 基礎上加入輻射一致性損失 $\mathcal{L}_{\text{rad}}$（$\alpha_1 > 0$），驗證陰影約束的獨立貢獻。
    \item \textbf{A5（GNN + $\mathcal{L}_{\text{rad}}$ + $\mathcal{L}_{\text{temp}}$）：}在 A4 基礎上加入時序平滑性損失 $\mathcal{L}_{\text{temp}}$（$\alpha_1, \alpha_2 > 0$），驗證時序約束的額外效益。
    \item \textbf{A6（完整 PI-ST-GNN）：}在 A5 基礎上加入風場阻滯損失 $\mathcal{L}_{\text{wind}}$（$\alpha_1, \alpha_2, \alpha_3 > 0$），為本研究之完整提案模型。
\end{itemize}


\subsection{消融實驗量化結果}
\label{subsec:ablation_results}

表~\ref{tab:ablation} 彙整五個消融變體與完整模型在測試集上的量化結果：

\begin{table}[htbp]
    \centering
    \caption[消融實驗量化結果比較]{消融實驗量化結果比較（真實場景測試集，166,188 預測點次）}
    \label{tab:ablation}
    \footnotesize
    \begin{tabular}{lccccp{2.8cm}}
        \toprule
        \textbf{模型變體} & \textbf{$R^2$} & \textbf{RMSE（°C）} & \textbf{MAE（°C）} & \textbf{準確率} & \textbf{組成說明} \\
        \midrule
        A0（MLP 基準線）    & 0.9926 & 0.455 & 0.230 & 97.4\% & 無 GNN，無物理約束 \\
        A3（GNN, 無物理）   & 0.9821 & 0.708 & 0.453 & 95.1\% & RGCN+LSTM \\
        A4（$+\mathcal{L}_{\text{rad}}$） & 0.9805 & 0.740 & 0.449 & 94.9\% & A3 + 輻射約束 \\
        A5（$+\mathcal{L}_{\text{temp}}$）& 0.9777 & 0.792 & 0.496 & 94.6\% & A4 + 時序約束 \\
        A6（完整模型）      & 0.9655 & 0.984 & 0.564 & 94.0\% & A5 + 風場約束 \\
        \midrule
        \textbf{最終（V5-300）} & \textbf{0.9753} & \textbf{0.833} & \textbf{0.524} & \textbf{93.9\%} & 完整 PI-ST-GNN，主要報告模型 \\
        \bottomrule
    \end{tabular}
\end{table}

\begin{figure}[htbp]
    \centering
    \includegraphics[width=0.85\textwidth]{../urban-thermal-gnn/04_training/figures/fig8_ablation.png}
    \caption[消融實驗各變體比較圖]{消融實驗各變體 $R^2$（左）與 RMSE（右）比較圖。}
    \label{fig:ablation}
\end{figure}


\subsection{消融結果討論：數值精度與物理可信賴性的取捨}
\label{subsec:ablation_discussion}

圖~\ref{fig:ablation} 以長條圖並列呈現各變體之 $R^2$（左）與 RMSE（右）。純 MLP 基準線（A0）於兩項數值指標上取得各變體中最佳表現（$R^2=0.9926$, RMSE$=0.455$°C）；此一數值優勢之訓練細節（樣本規模、空間平滑效應等）屬模型工程層面之技術分析，完整說明見附錄~\ref{appx:ablation_numeric_detail}。就本研究之設計輔助宗旨而言，更關鍵的觀察在於：A0 缺乏物理約束，其預測可能違反輻射陰影之空間一致性（即輸出「陰影區比日照區更熱」之物理悖論），此類違反難以被 $R^2$ 或 RMSE 這類彙總指標捕捉，卻直接構成建築設計決策之誤導風險。

從 A3 至 A6 逐步加入物理損失，RMSE 呈現單調上升（A3 $0.708\to$A4 $0.740\to$A5 $0.792\to$A6 $0.984$°C），但每項約束各自針對特定的熱力學一致性問題發揮作用：$\mathcal{L}_{\text{rad}}$ 之加入（A3$\to$A4）強制模型學習陰影梯度之方向性，改善局部輻射一致性；$\mathcal{L}_{\text{temp}}$ 之加入（A4$\to$A5）防止逐時預測出現非物理跳躍，改善時序軌跡的熱慣性自洽性；$\mathcal{L}_{\text{wind}}$ 之加入（A5$\to$A6）約束高層建築背風側之熱累積極端值，改善風場阻滯情境的物理合理性。這些一致性提升難以被純數值指標捕捉，卻是建築師採信模型輸出、將其用於實際設計決策之前提。換言之，完整 PI-ST-GNN 以較程序化資料更明顯的數值精度取捨，換取設計輔助工具最不可或缺的特質：預測結果在物理上站得住腳。

須說明者：本節五個消融變體之訓練，係於真實場景資料集（V5-300 相同切分，動態邊快取已修正，\texttt{shadow}／\texttt{veg\_et}／\texttt{convective} 三類關係均透過真實邊傳遞訊息）上重新執行，故其異質圖訊息傳遞已包含全部五類關係型邊緣，不受 V1--V4 期間動態邊缺陷之影響（該缺陷之完整記錄見附錄~\ref{appx:dynamic_edges_fix}）。五個變體彼此僅在物理損失項組合上有差異，圖結構部分完全一致，故變體間之相對比較仍屬內部一致；本節 A6 變體與表~\ref{tab:main_results}之 V5-300 官方主要報告模型間之數值差異說明，見附錄~\ref{appx:ablation_numeric_detail}。


\section{推理效能對互動式設計工作流之意義}
\label{sec:performance}

\subsection{推理延遲基準測試}
\label{subsec:inference_timing}

本研究對 PI-ST-GNN 進行 200 次重複推理計時（單場景，程序化主線之固定 6,241 個空氣節點網格，11 時步），結果如表~\ref{tab:timing} 所示；此一節點數為程序化主線之固定取樣網格，V5-300 所用之真實場景節點數（136--400，第~\ref{subsec:overall_metrics}節）較此更小，故下列延遲數字可視為 V5-300 實際推理延遲之保守上界，而非低估：

\begin{table}[htbp]
    \centering
    \caption[推理延遲基準測試]{GNN 推理延遲基準測試（RTX 5070 Ti, 200 次重複）}
    \label{tab:timing}
    \small
    \begin{tabular}{lccc}
        \toprule
        \textbf{評估方法} & \textbf{平均延遲（ms）} & \textbf{最小延遲（ms）} & \textbf{最大延遲（ms）} \\
        \midrule
        PI-ST-GNN（本研究） & \textbf{61.5 ± 24.4} & 37.1 & 185.6 \\
        傳統 CFD 模擬（參考值）& $\sim$2--8 小時 & --- & --- \\
        Ladybug Tools 微氣候模擬 & 數分鐘至數十分鐘 & --- & --- \\
        \bottomrule
    \end{tabular}
\end{table}

此處關注的重點並非推理延遲的絕對數字，而在於代理模型的評估時間尺度是否落在設計迭代與最佳化搜尋可承受的範圍內。相較於傳統 CFD（單次求解約 2--8 小時）與 Ladybug Tools 微氣候模擬（數分鐘至數十分鐘）需以離線批次方式執行，代理模型將單次適應度求值壓縮至可即時互動的量級，使原本受限於模擬求值成本的兩種工作流程得以成立：其一，設計師於 Rhino／Grasshopper 環境中調整建築量體後，可在設計操作的時間尺度內取得全場域更新的 UTCI 熱圖；其二，將代理模型嵌入 NSGA-II 之適應度函數，使數百代、每代數十至數百個體之最佳化搜尋得以在可行的運算時程內完成，而非傳統模擬驅動流程所需之數週時間。本研究之貢獻在於方法工作流程之整合可行性，而非單純追求運算速度之極大化。


\subsection{訓練資料集覆蓋範圍驗證}
\label{subsec:coverage}

為驗證 300 個真實訓練場景對都市形態空間的覆蓋充分性，本研究分析訓練集在主要幾何特徵（FAR、BCR、建築棟數、樓高）上的統計分佈。由於場景直接取自真實 OSM 建物量體而非法規約束下之程序化生成，其分布較程序化幾何遠為寬廣且不均勻：FAR 中間 90\% 範圍（第 5--95 百分位）為 0.54--12.53（全距 0.23--40.83，涵蓋自低密度郊區至高密度都市核心之真實建物量體），BCR 中間 90\% 範圍為 0.18--1.00（全距 0.08--1.00），每場景建物棟數 3--208 棟（平均 11.3 棟），建物簷高 3--162 公尺（平均 12.6 公尺）。此一遠較程序化幾何寬廣之真實形態覆蓋範圍，是第~\ref{chapter:results}章以真實都市建物量體進行案例驗證、以及第~\ref{sec:nsga2}節 NSGA-II 最佳化搜尋時，判斷結果是否落入分布外推（Out-of-Distribution）風險的重要參照基準——第~\ref{subsec:pareto_morphology}節之最佳化結果所落之 FAR 區間即落於此訓練分布之中間範圍內。


\section{真實場景資料集：版本演化與模擬引擎驗證}
\label{sec:v4_validation}

本研究之真實場景資料集與模擬引擎歷經五個版本之演化，於此簡述以說明實驗歷程與本研究最終採用 V5-300 為主要報告模型之原因。V1 以八維空氣節點特徵於程序化幾何上建立基礎預測能力；V2 增列地表溫度為第九維特徵；V3 進一步納入街谷高寬比（$H/W$）為第十維特徵；上述三版皆以程序化蒙地卡羅幾何與典型氣象年為訓練資料。V4 為資料版本之根本轉變——由程序化幾何全面改採真實 OSM 建物、ETH 樹冠幾何，並以真實環境部 IoT 測站之逐時實測氣溫濕度驅動模擬（場址、建物、樹冠與 IoT 氣象等地理資料之完整說明見第~\ref{subsec:v4_real_dataset}節），惟其微氣候模擬仍採自製 Python 物理近似引擎。V5 在完全承襲 V4 真實地理資料與氣象邊界條件之基礎上，將模擬引擎全面替換為 Radiance／EnergyPlus 官方模擬流程（第~\ref{sec:ladybug_workflow}節），並同時修正下述動態邊缺陷，是本研究最終採用之主要報告模型，其完整訓練與測試結果已於第~\ref{sec:evaluation}節（表~\ref{tab:main_results}、圖~\ref{fig:training_curves}--\ref{fig:spatial_comparison}）報告，本節不再重複列出。

須揭露者：V1--V4 訓練時所用之異質圖動態邊快取（\texttt{shadow}／\texttt{veg\_et}／\texttt{convective} 三類關係）與程序化主線曾發生之缺陷同屬一類——快取檔案從未於真實場景資料集上被實際建立，故 V1--V4 之異質圖訊息傳遞實際上僅透過 \texttt{semantic} 與 \texttt{contiguity} 兩類靜態邊進行；V4 當時表面報告之測試集數值（$R^2 \approx 0.99$ 量級）即在此一未修正狀態下取得，不具方法完整性，不應作為有效驗證結果引用。此一缺陷直到 V5-300 以 \texttt{16\_build\_dynamic\_edge\_cache\_v5.py}（真實逐場景太陽方位角、真實逐月中央氣象署風向）重新建構動態邊快取後方才修正，完整技術記錄見附錄~\ref{appx:dynamic_edges_fix}。

圖~\ref{fig:v5_comparison} 就數個未見過之真實測試場景，並列呈現其真實 OSM 幾何、Radiance／EnergyPlus 模擬之真值 UTCI 場、V5-300 模型預測場與逐點誤差，作為表~\ref{tab:main_results}量化指標之空間視覺化補充。

\begin{figure}[htbp]
    \centering
    \includegraphics[width=\textwidth]{../urban-thermal-gnn/04_training/figures/fig_v5_scene_comparison.png}
    \caption[真實測試場景之預測比較]{以真實地理環境與 Radiance／EnergyPlus 引擎訓練之 PI-ST-GNN（V5-300）於未見過真實測試場景之預測比較（14:00）：由左至右為真實 OSM 幾何、真值 UTCI（Radiance／EnergyPlus 模擬）、模型預測、逐點誤差。幾何欄之虛線框標示 $80\times80$ 公尺基地邊界，實色為基地內建物與喬木，淡色為基地外仍參與遮蔭運算之周邊紋理脈絡。}
    \label{fig:v5_comparison}
\end{figure}

須對照說明者：V5-300 之 $R^2 = 0.9753$ 略低於本研究另一條程序化主線（第~\ref{sec:evaluation}節前述、以合成幾何訓練）之 $0.9956$，此差異為合理且可預期——Radiance／EnergyPlus 模擬之空間變異本質上較自製近似式更為劇烈與寫實（局部 MRT 極端值可達 77\,$^\circ$C 以上），且建物外殼構造採 ASHRAE 標準假設而非逐棟實測材料，真實 OSM 幾何與真實 IoT 觀測亦含更高之雜訊與異質性（如樓高標籤稀疏、單一氣象站空間代表性等）。然而模型於\textbf{未見過之真實場景}上仍維持 RMSE 僅 $0.83\,^\circ\text{C}$、遠小於 UTCI 熱壓力等級間距，此為模型於真實地理環境與物理模擬雙重條件下之外部效度直接證據，亦是本研究方法論上最嚴謹之驗證結果，故雖數值指標略低於程序化主線，仍以 V5-300 作為全文主要報告模型。


\section{小結}
\label{sec:chapter4_summary}

本章以測試集量化評估、消融實驗與推理效能基準測試，驗證了 PI-ST-GNN 作為都市微氣候代理模型之有效性與可信賴性。本研究最終採用之主要報告模型（V5-300：真實地理資料 $+$ Radiance／EnergyPlus 模擬引擎 $+$ 已修正動態邊快取，附錄~\ref{appx:dynamic_edges_fix}）在 166,188 個真實場景測試樣本上達到 $R^2=0.9753$、RMSE$=0.833$°C、熱壓力等級分類準確率 93.9\%，為本研究方法論上最嚴謹之驗證結果（第~\ref{sec:v4_validation}節）；消融實驗進一步揭示，純數值指標上表現最佳的 MLP 基準線，缺乏圖結構與物理約束所提供的空間因果建模能力與預測物理一致性，而完整 PI-ST-GNN 以輕微數值犧牲換取了設計輔助工具不可或缺的熱力學可信賴性；代理模型將單次適應度求值壓縮至可即時互動之時間量級，則為第~\ref{chapter:results}章之多目標最佳化與熱舒適人行動線評估，提供了得以整合入設計操作流程之核心運算引擎。
